\subsection{The middle unipotent stabilizer}
\label{sec:unipotent-stabilizer-covolume}

The weight of a flag class is a ratio of stabilizer covolumes.
Its dependence on the large definite rank comes from one
unipotent quotient.  We prove the precise estimate for that
quotient.  The exponent of the integral lattice defect is
the definite rank itself; no factor depending on the degree
of the field is lost.

\paragraph{The normal form and its lattices.}
Let $E/F$ be a CM extension unramified at every finite place,
and put $\mathcal O=\mathcal O_E$.  In this subsection the
Hermitian form $h$ is linear in the first variable, as in the
source's normal form.  Replacing $h(x,y)$ by $h(y,x)$ converts
to the opposite convention without changing any lattice,
index or covolume below.

Let $0\leq j\leq r$.  Write
\[
 X=\langle e_0,\ldots,e_{r-1}\rangle_E,\qquad
 X'=\langle e'_0,\ldots,e'_{r-1}\rangle_E
\]
for two totally isotropic spaces.  Suppose the pairing between
their first $j$ vectors is nondegenerate, and that
$e'_a=c e_a$ for $j\leq a<r$, where $c\bar c=1$.
Put $X_j=\langle e_a:a<j\rangle_E$ and
\[
 I=X\cap X'=\langle e_j,\ldots,e_{r-1}\rangle_E.
\]
Let $z$ be perpendicular to $X+X'$, with $h(z,z)=-T\ne0$.
Choose a totally isotropic dual space with basis
$f_j,\ldots,f_{r-1}$ satisfying
\[
 h(e_b,f_a)=\delta_{ba}\quad(0\leq b<r),\qquad
 h(e'_b,f_a)=0\quad(b<j),\qquad h(z,f_a)=0.
\]
Let $W$ be the orthogonal complement of all these vectors,
and write $k=\dim_E W$.  Thus
\[
 V=\bigl(X_j+\langle e'_a:a<j\rangle_E\bigr)
       \mathbin\perp(I+\langle f_a:a\geq j\rangle_E)
       \mathbin\perp Ez\mathbin\perp W.
 \tag{US1}
\]
Set $V_0=X+X'+Ez$.  Then $V_0^\perp=W+I$.
The source's normal form has exactly these spaces and vectors.
\cite{unipotent-stabilizer-source}

Fix nonzero fractional ideals $A_0,\ldots,A_{r-1}$ invariant
under conjugation.  In the original construction
$A_a=\mathfrak d^{-(r-a)}$.  Let $\Lambda$ be an integral
full $\mathcal O$-lattice.  Its prescribed flag condition is
\[
 \{t\in E:t e_a\in\Lambda+X_a\}=A_a\quad(a<r),
 \qquad \{t\in E:t z\in\Lambda+X\}=\mathcal O,
 \tag{US2}
\]
where $X_a=\langle e_b:b<a\rangle_E$.
This is the condition on successive flag quotients used for
the adapted normal form.  The original row flags satisfy the
stronger literal weighted-lattice containment used in
Section~\ref{shared-integral-leading-minors}; no substitution
of those two interfaces is being made.

Define full lattices
\[
 M=\Lambda\cap W,\qquad L=(\Lambda+X)\cap W,\qquad
 J=(\Lambda+X)\cap(W+Ez),
\]
and the reference lattice
\begin{align*}
 P_0&=\sum_{a<r} A_a e_a+\sum_{a<j} A_a e'_a
                      +\mathcal O z+\sum_{a=j}^{r-1}A_a^{-1}f_a,\\
 \Lambda_{\mathrm{ref}}&=P_0+J.
 \tag{US3}
\end{align*}
These are the source's definite intersection, reference
module and reference lattice.  The redundant perpendicular
conditions in its definitions disappear because
$W+Ez\subset(X+X')^\perp$.  Equation (US2) gives
$J\cap Ez=\mathcal O z$.  Consequently projection onto $W$
in (US1) sends
\[
 \Lambda_{\mathrm{ref}}\cap(W+I)\quad\hbox{onto }L.
 \tag{US4}
\]
Indeed the $z$-coordinate of the $P_0$ term lies in
$\mathcal O z$, so any cancellation of that coordinate
can be performed inside $J$.

Put
\[
 b=[\Lambda\cap X:(\Lambda\cap X_j)+(\Lambda\cap I)],
 \qquad d=[L:M].
 \tag{US5}
\]
Both are positive integers.  All indices mean cardinalities
of additive quotients, including the residue degrees.

\paragraph{The group and the meaning of the weight.}
Let $U$ be the group of unitary transformations fixing $V_0$
pointwise and inducing the identity on $V_0^\perp/I$.
Let $Z\subset U$ be the subgroup fixing $W$ pointwise.
This is precisely the source's group $R_1$ and its subgroup
$Z_I$.  Indeed its flag unipotent group consists of the
unitary transformations acting identically on every
successive quotient of $X_0\subset\cdots\subset X_r=X$
and on $X^\perp/X$.  The two such flag groups are
subgroups, so their intersection is already the subgroup
generated in the definition of $R$.  An element of $R$
fixing $V_0$ sends $w\in W$ to $w$ modulo both $X$ and
$X'$, hence modulo $I$.  Conversely an element of $U$
fixes all flag vectors and acts identically modulo $X$
on $X^\perp=X+Ez+W$; the same argument applies to
$X'$.  It therefore belongs to both original flag groups.
This proves the asserted equality of groups.
For a full lattice $B$ write
$\Gamma_B=\{u\in U:uB=B\}$.  For commensurable subgroups
$A,B$ define
\[
 \operatorname{ivol}(A,B)=
            \frac{[A:A\cap B]}{[B:A\cap B]}.
 \tag{US6}
\]
This definition is the source's inverse covolume ratio.  It
introduces no choice of Haar measure.  The groups
$(\Gamma_{\Lambda_{\mathrm{ref}}})Z$ and
$(\Gamma_\Lambda)Z$ have the common normal subgroup $Z$;
their ratio is computed in the additive quotient $U/Z$.

\begin{theorem}[The middle stabilizer estimate]
\label{thm:middle-unipotent-stabilizer}
With the preceding data,
\[
 \operatorname{ivol}
       (\Gamma_{\Lambda_{\mathrm{ref}}}Z,\Gamma_\Lambda Z)
                         \leq b^k d^{\,r-j}.
 \tag{US7}
\]
For $j=r$, both sides are $1$.  In the original split normal
form, $d\leq xi$, where $x$ and $i$ are its two successive
fixed-lattice indices.  Thus (US7) supplies the source's exact
bound
\[
                         b^k(xi)^{r-j}.
 \tag{US8}
\]
\end{theorem}

\begin{proof}
First, there is $w_0\in\mathcal O$ with $w_0+\bar w_0=1$.
To see this, localize the trace ideal at a prime of
$\mathcal O_F$.  In the split case the trace is addition.
In the unramified quadratic case its reduction is the
surjective trace of finite fields; Nakayama's lemma gives
surjectivity over the local ring.  The trace ideal is
therefore $\mathcal O_F$, which proves the assertion
globally.  This includes primes above $2$.

For $a\geq j$ define
\[
             \rho_a(u)=\operatorname{pr}_W(uf_a-f_a).
\]
Since $u$ fixes $V_0$, the displacement lies in $W+I$.
The action on that space modulo $I$ is the identity, so
$\rho=(\rho_a)_{a\geq j}$ is a homomorphism
$U\to W^{r-j}$, with the additive group law on the target.
Its kernel is $Z$.  For example, if every row is zero,
$uf_a-f_a\in I$.  Pairing with $uw$ for $w\in W$, and
using $uw-w\in I$, shows that every coordinate of
$uw-w$ on the basis $e_j,\ldots,e_{r-1}$ is zero.
The converse follows by the same pairing.

The images of both lattice stabilizers under $\rho$ are
full additive lattices.  An upper lattice bound follows by
applying their elements to a fixed set of rational basis
vectors with bounded denominators.  For the lower bound,
the explicit transformations below, with the isotropic and
definite vectors multiplied by sufficiently divisible
integers, produce full lattices in every row.  Such multiples
can be chosen in both $B$ and its Hermitian dual for any
full lattice $B$.  The formulas and their inverses then
preserve $B$.  This also proves the commensurability required
in (US6).

For $u\in\Gamma_{\Lambda_{\mathrm{ref}}}$ and
$t\in A_a^{-1}$, the vector $tf_a$ lies in $P_0$.
Hence $t(uf_a-f_a)$ lies in the left side of (US4).
Projection to $W$ gives $t\rho_a(u)\in L$.  Invertibility
of $A_a$ implies
\[
 \rho(\Gamma_{\Lambda_{\mathrm{ref}}})
                         \subset\prod_{a=j}^{r-1}A_aL.
 \tag{US9}
\]
No freeness assumption on $L$ or $A_a$ is used.

For each $a\geq j$, let $C_a$ be the ideal of coefficients
$h(f_a,x)$ of vectors $x\in\Lambda\cap I$ satisfying
$h(f_b,x)=0$ for $b>a$.  Since conjugation preserves $A_a$,
the leading coefficient condition (US2) gives $C_a\subset A_a$.
The quotient of $\Lambda\cap X$ by $\Lambda\cap X_j$ is
the projected lattice in $I$, with successive leading
ideals $A_j,\ldots,A_{r-1}$.  The corresponding leading
ideals of $\Lambda\cap I$ are the conjugates of $C_a$.
The index formula for these filtrations gives
\[
                     \prod_{a=j}^{r-1}[A_a:C_a]=b.
 \tag{US10}
\]
One may check this formula after localization, where it is
the determinant of a triangular change of lattice basis;
the finite local quotients multiply to the global index.

Here is the required transformation inside the actual
integral stabilizer.  For $x\in\Lambda\cap I$ and
$v\in M$, set $u=-v$ and $q=-w_0h(u,u)$.  The map
\[
 E_{x,u,q}(y)=y+h(y,u)x-h(y,x)u+q h(y,x)x
 \tag{US11}
\]
is unitary because $h(x,x)=h(x,u)=0$ and
$q+\bar q=-h(u,u)$.  Expanding its Hermitian pairing
verifies this identity directly.  Its inverse is
$E_{x,-u,\bar q}$.  Integrality of $\Lambda$ implies
that both formulas preserve $\Lambda$.
They fix $V_0$ and act trivially on $W$ modulo $I$, so
$E_{x,u,q}\in\Gamma_\Lambda$.  Its rows are
\[
                     \rho_a(E_{x,-v,q})=h(f_a,x)v.
 \tag{US12}
\]
The signs in (US12) follow by projecting (US11) to $W$;
all terms proportional to $x$ disappear.

Order the rows from the largest index downwards.  Let $Q_a$
be the projection to row $a$ of those elements of
$\rho(\Gamma_\Lambda)$ whose rows with index greater than
$a$ vanish.  Equation (US12), and addition in the quotient
$U/Z$, give
\[
                              C_aM\subset Q_a.
 \tag{US13}
\]
Use product Lebesgue measure in any fixed real coordinates
on $W^{r-j}$.  Covolumes multiply along the row filtration.
Writing $A=\rho(\Gamma_{\Lambda_{\mathrm{ref}}})$ and
$B=\rho(\Gamma_\Lambda)$, (US9) and (US13) give
\[
 \frac{\operatorname{covol}B}{\operatorname{covol}A}
 \leq\prod_{a=j}^{r-1}
             \frac{\operatorname{covol}(C_aM)}
                  {\operatorname{covol}(A_aL)}
 =\prod_{a=j}^{r-1}[A_aL:C_aM].
\]
For any invertible ideals $C\subset A$ and full lattices
$M\subset L$ of rank $k$,
\[
                       [AL:CM]=[A:C]^k[L:M].
 \tag{US14}
\]
This follows locally by using a basis over the discrete
valuation ring, and then multiplying the resulting finite
quotients.  Equations (US10) and (US14) prove (US7), because
the covolume ratio is precisely (US6) after quotienting by
the common kernel $Z$.

For the original split normal form,
$\Lambda\cap(W+I)=M+(\Lambda\cap I)$.  The definite
quotient lattice is $L$, and its fixed lattice $K$ satisfies
$M\subset K\subset L$.  In the source notation
$x=[L:K]$ and $i=[K:M]$ after projection modulo $I$.
Thus $d=xi$; the weaker inequality $d\leq xi$ used by
the source follows as well.  More explicitly, $K$ is the
set of $w\in L\cap L'$ admitting lifts
$\lambda-w\in X$ and $\lambda'-w\in X'$, with
$\lambda,\lambda'\in\Lambda$ and
$\lambda-\lambda'\in g\Lambda+(\Lambda\cap X)$;
here $L'=(\Lambda+X')\cap W$.
This is its original diagonal congruence condition.
\end{proof}

The bound proves the large-rank part of the original
stabilizer estimate.  To finish the complete weighted
flag-class sum, one must still construct the common
unipotent transport placing the rank-$(2r+1)$ lattice
between the specified ideal multiples of $P_0$, and
verify the top and central stabilizer bounds on those
same integral lattices.  Those obligations are not
consequences of (US7).
